\chapter{来源与读法索引}
\label{app:sources}

\phantomsection
\section*{来源索引}
\label{app:source-index}

本索引列出本卷账本和正文实际依赖的材料路径，并说明每一类材料能承担、不能承担什么判断。引文应回到本索引所列的版本、卷次、日期、档号或版面；电子全文库只用于检索。事件锚点的最小证据链另见账本，故这里不把一条诏令、一个战报或一份条约夸大为对地方社会的全知叙述。要从正文的事件入口回查，先看\hyperref[app:event-navigation]{事件导航}；要按时期、行动空间或图件重走叙事，再看\hyperref[app:index]{读者索引}。下面每一行都把材料类型、可问的问题、在本卷的用法和判断边界并列，而不是把书名当作结论。

\section{宫廷、编年与制度}

\begin{tabular}{@{}p{0.23\textwidth}@{}p{0.20\textwidth}@{}p{0.25\textwidth}@{}p{0.22\textwidth}@{}}
\textbf{材料、作者／编者与版本类型} & \textbf{时段或问题} & \textbf{本卷用法} & \textbf{不能承担的判断}\\
\hline
\textbf{材料、作者／编者与版本类型} & \textbf{时段或问题} & \textbf{本卷用法} & \textbf{不能承担的判断}\\
\hline
《清实录》（清廷纂修；中华书局影印本，按各朝实录卷日引用） & 1644--1908 的上谕、任免、军政与奏报 & 建立清廷获报、命令和官方日期的编年骨架 & 不证明命令已执行；军功、伤亡、人口和对手动机须另证。\\
\hline
《宣统政纪》（清末国史馆材料；中华书局影印本） & 1908--1912 的朝廷程序与宣统朝政务 & 同内阁、军机和报刊材料校对退位前后的程序 & 不能把公告日等同各省、蒙古或边地的政治转向。\\
\hline
赵尔巽等编《清史稿》（中华书局点校本） & 全期的制度、人物、地理和邦交检索 & 作为后出索引，追查实录、档案与地方材料的线索 & 未定稿的民国叙述不能单独证明数字、边界或褒贬。\\
\hline
王先谦等编《东华录》、国史馆起居注和《清史列传》（中华书局影印本或可靠影印本） & 宫廷政务、诏令摘要和人物叙述 & 对照实录的异文与国史叙述来源 & 摘编、抄录和后出的评语不可当作原件日期。\\
\hline
中国第一历史档案馆编《清宫朱批奏折档案总汇》及宫中朱批奏折原件／影印件 & 康熙以后的人事、灾荒、财政与军务 & 分开奏者的呈报、皇帝朱批和后续处分 & 一份奏折只证明一位官员的陈述与批答，不证明地方结果。\\
\hline
中国第一历史档案馆藏军机处档、随手登记档、议复档与《军机处奏折录副》 & 雍正以后紧急军政、边务、外交与调兵 & 追踪案件在中央的流转及跨省协调 & 附件、口谕、地方接收和实际调度常不完整。\\
\hline
中国第一历史档案馆藏内阁大库题本、部院档与《明清档案》影印整理 & 赋役、河工、司法、工程和常规行政 & 识别题报、部议、题准与核销等程序节点 & 题准、定额或报销不等于已经征到、修成或发到。\\
\hline
《大清会典》《大清会典事例》《大清律例》（中华书局或国家图书馆影印本，按修纂年使用） & 官制、律例、财政名目和制度变化 & 界定官名、规范与制度语汇 & 定制会压平例外和地方变通，须由案件、账册与档案检验。\\
\hline
\end{tabular}

\section{战争、边疆与跨境材料}

\begin{tabular}{@{}p{0.23\textwidth}@{}p{0.20\textwidth}@{}p{0.25\textwidth}@{}p{0.22\textwidth}@{}}
\textbf{材料、作者／编者与版本类型} & \textbf{时段或问题} & \textbf{本卷用法} & \textbf{不能承担的判断}\\
\hline
\textbf{材料、作者／编者与版本类型} & \textbf{时段或问题} & \textbf{本卷用法} & \textbf{不能承担的判断}\\
\hline
《平定三逆方略》《平定准噶尔方略》《平定回部方略》（清廷纂修；文渊阁《四库全书》影印本，并校军机档） & 三藩、蒙古、准噶尔与回部战争 & 追踪清廷如何编排调兵、奖惩和凯旋 & ``平定''、兵数、歼获与疆界是国家战时语言，须和地方及多语材料互校。\\
\hline
理藩院满汉蒙藏档、盟旗档和满文《清实录》（第一历史档案馆整理本或影印件；按档号、语种引用） & 蒙古、青海、西藏与北部边疆 & 说明册封、牧地、寺院、驻防和文书分类的制度位置 & 清廷分类不是唯一政治视角；译名、地名与盟旗关系须保留不确定性。\\
\hline
藏文寺院文书、碑刻与传记的可靠译注本；驻藏大臣档 & 西藏、青海、川西和喜马拉雅交通 & 将寺院、噶厦、驻防和贸易放入各自时间线 & 不可用汉文驻藏叙述代替藏文地方政治，反之亦然。\\
\hline
维吾尔／察合台文文书的可靠释读本，连同伊犁将军档和新疆地方档 & 伊犁、天山南北、绿洲水利与贸易 & 检验征服后屯田、税役、土地和跨境网络的局部变化 & 文书出土地、保存史和译读有限，不能外推全新疆或固定现代族称。\\
\hline
《尼布楚条约》满、拉丁、俄文各本及俄国使团、边防档案的校勘／影印本 & 黑龙江、阿尔巴津和多语边界交涉 & 对照条文、翻译和签署程序 & 不能从河名和条款投射出精确现代边界或居民的实际知晓。\\
\hline
《筹办夷务始末》、总理衙门档、中外条约各文与换约文本 & 鸦片战争以后外交、口岸、边界与赔款 & 分开照会、签署、批准、换约和执行的时间 & 外交辞令、单一译本或条约权利均不等于实际控制。\\
\hline
英国、俄国、法国、日本等国的外交／军政／领事档案（原档、官方档案刊本或可核影印本）；《朝鲜王朝实录》《承政院日记》《大南实录》及琉球／日本、俄蒙中亚文献（原文影印或有底本的译注本，按语种引用） & 战争、海防、朝贡、边务与相邻政权的行动 & 将俄罗斯、朝鲜、越南、日本及边地行动者置入各自的文书链，并与清方材料校读 & 外交、领事与海军报告服务本国情报和帝国利益；礼仪称谓、译名与观察地点不同，不能视作中立观察或直接等同主权。\\
\hline
\end{tabular}

\section{地方社会、海疆与近代公共信息}

\begin{tabular}{@{}p{0.23\textwidth}@{}p{0.20\textwidth}@{}p{0.25\textwidth}@{}p{0.22\textwidth}@{}}
\textbf{材料、作者／编者与版本类型} & \textbf{时段或问题} & \textbf{本卷用法} & \textbf{不能承担的判断}\\
\hline
\textbf{材料、作者／编者与版本类型} & \textbf{时段或问题} & \textbf{本卷用法} & \textbf{不能承担的判断}\\
\hline
府州县志、边防志与台湾府县志（据各书修纂年、版本和卷次引用） & 城池、地名、学校、水利、港口、灾异和地方记忆 & 定位城市、河口、道路与地方制度的长期变化 & 后修、沿袭与地方自我呈现显著；一部志书不可证明实时事件。\\
\hline
巴县档案及省、府、州县案件、赋役和契约档（四川省档案馆等整理／影印本，按档号） & 诉讼、地租、借贷、差役与基层行政 & 将地方财政和社会选择放回当事人、地点和程序 & 保存地域高度不均，一个案件或一县不能代表全国。\\
\hline
账簿、族谱、会馆／行会文书、碑刻和墓志（原件、拓本或可靠著录本） & 商业、迁徙、宗族、慈善、土地与信贷 & 限定一个组织、市场或家庭的资源关系 & 保存偏向识字与有产者；自称、追赠和登记均须外部核查。\\
\hline
台湾档案、契约、港口记录、航海日志与荷兰／西班牙／日文材料的可靠影印／译注本 & 郑氏至清治台湾、海峡贸易、原住民与海防 & 区分海禁命令、航行、港口贸易、设治和地方接触 & 殖民、官府和商人材料各有土地与主权诉求；原住民自述保存稀少。\\
\hline
《海关贸易报告》《中国海关十年报告》及总税务司档案（海关总税务司刊本／影印本） & 1859年后口岸贸易、关税、航运和城市信息 & 使用有口岸、货类、年度和机构口径的统计 & 覆盖制度内口岸贸易，不等于内地、走私或全国经济总量。\\
\hline
《申报》、京报、公报及上海、天津、汉口、广州、香港报刊影印／数据库检索件 & 清末战事、改革、革命、商业与城市公共信息 & 区分消息首发、转载和公告进入公共流通的时间 & 报刊的审查、通讯延误和读者市场限制其代表性。\\
\hline
日记、书信、回忆录与传教团体档案的可靠校点／影印本 & 个体的行程、感知、赈济与信息网络 & 为具体人、地、日序提供局部证词和待查线索 & 后改、删节、立场和后见之明明显；不能替未留文字者发言。\\
\hline
\end{tabular}

\section{按问题回查的材料路线}

下列路线对应\hyperref[app:event-index]{按时期与区域阅读}和\hyperref[app:polity-index]{行动空间与主题}中的问题入口。它们不是为单一事件预先裁定证据，而是让读者知道应当把哪一种文书、哪一种版本与哪一种局部经验放在一起读；具体引文仍须落到卷日、档号、日期或版面。

\begin{description}
  \item[入关、南方征服、海峡与草原（1644--1691）] 先以《清实录》的卷日和内阁大库题本确立清廷收到什么消息、何时作何命令；三藩与海峡战事再并读《平定三逆方略》、福建及台湾府县志、港口记录和航海日志；尼布楚与多伦诺尔则加读《尼布楚条约》的满、拉丁、俄文各本、俄方档案及理藩院满汉蒙藏档。它们在本卷中用来拆开北京决策、航路或草原中介与地方接收，不能把战报、条约或册封文字改写成征服完成、精确边界或居民认同。
  \item[财政、军机与西北／高原的治理（1723--1797）] 雍正财政和军机房的程序以宫中朱批奏折、军机处档、《大清会典》及《大清会典事例》追读；伊犁、乌什、金川和青藏问题须把《平定准噶尔方略》《平定回部方略》与理藩院档、伊犁将军档、满蒙藏文或维吾尔／察合台文文书的可靠释读并置。本卷据此区分呈报、批答、调兵与驻防后的地方条件，不能由定制、凯旋叙事或一语种材料推出税役实际、伤亡、牧地关系或全区域秩序。
  \item[口岸、条约与战争（1839--1864）] 虎门、南京、北京与各口岸的程序从《筹办夷务始末》、总理衙门档、条约中外文签署本及批准／换约文本入手，并与英、法等外交或军政档、中文及外文报刊对读。它们在本卷中分别追踪照会、签署、公告和信息流通，不能以一份译本、一次刊载或一方战报证明条约已经执行，或概括港口与内地社会的经历。
  \item[地方社会、灾荒、重建与海疆] 地方志、巴县及省府州县案件、赋役和契约档，连同账簿、会馆文书、碑刻、墓志和日记，构成本卷追问地租、诉讼、迁徙、赈济及地方组织的入口；台湾问题另并读台湾档案、契约、港口记录和荷兰／西班牙／日文材料的可靠影印或译注本。这些材料只让读者回到一地、一案或一段网络的时间与关系，不能把保存较多的县、城市、识字者或殖民记录外推为全国、全岛或未留文字者的经验。
  \item[晚清改革、公共信息与帝国终结（1877--1912）] 赈济、伊犁、甲午、新政和革命首先分辨《清实录》末期卷日与《宣统政纪》、宫中奏折和军机／总理衙门档的文件层次；再以《申报》、京报、公报及各地中文、外文报刊的影印件核查首发、转载和公告，以《海关贸易报告》《中国海关十年报告》限定口岸统计，必要时与社团刊物、日记、书信和地方档案互校。它们在本卷中把宣言、报导、统计、地方变化和退位程序分开，不能从报刊声量、海关数值、政治主张或退位诏书直接推出全国民意、经济总量、军政控制或同步的政治转换。
\end{description}

\section{书系读法入口}

二十四史是书系的长时段正史骨架：司马迁《史记》、班固《汉书》、范晔《后汉书》、陈寿《三国志》、房玄龄等《晋书》、沈约《宋书》、萧子显《南齐书》、姚思廉《梁书》《陈书》、魏收《魏书》、李百药《北齐书》、令狐德棻等《周书》、魏徵等《隋书》、李延寿《南史》《北史》、刘昫等《旧唐书》、欧阳修、宋祁《新唐书》、薛居正等《旧五代史》、欧阳修《新五代史》、脱脱等《宋史》《辽史》《金史》《元史》《明史》。本卷以中华书局点校本这一通用阅读版本，追问清以前制度、边地和区域历史的承接；它们不是清代逐年事实的证据。清代的年序与制度问题须回到《清实录》、清末档案、地方文献和多语材料，不能把《清史稿》误列入二十四史。

钱穆《国史大纲》（商务印书馆再版本）、钱穆《历代政治得失》（生活·读书·新知三联书店再版本）、张荫麟《中国史纲》（中华书局再版本）和吕思勉的通史、断代史著作（上海古籍出版社或中华书局再版本）供读者形成时代大势、制度得失与社会经济长线的问题意识。本卷不以它们单独确证逐年日期、精确数字或人物动机。《古文观止》（吴楚材、吴调侯编；中华书局注释本）仅作为古文阅读入口；遇到清代诏令、条约或地方文字，必须回到原始文本、可靠底本及其形成语境。
